% Mathematics
%%%%%%%%%%%%%%%%%%%%%%%%%%%%%%%%%%%%%%%%%%%%%%%%%%%%
% ریاضیات
%%%%%%%%%%%%%%%%%%%%%%%%%%%%%%%%%%%%%%%%%%%%%%%%%%%%
\usepackage{amsmath}
\usepackage{amssymb}
\usepackage{amsthm}
\usepackage{mathtools}
\usepackage{makeidx}
\usepackage{booktabs}
%%%%%%%%%%%%%%%%%%%%%%%%%%%%%%%%%%%%%%%%%%%%%%%%%%%%
% اعداد و واحدها
%%%%%%%%%%%%%%%%%%%%%%%%%%%%%%%%%%%%%%%%%%%%%%%%%%%%
\usepackage{siunitx}

\sisetup{
output-decimal-marker={/},
group-separator={,},
group-minimum-digits=4
}
%برای قرار دادن کاما برای جدا کردن ارقام (جداکننده هزارگان) در اعداد فارسی (یا انگلیسی) 
% در لاتک و زی‌پرشین (XePersian)،
% بهترین روش استفاده از بسته
% siunitx
% است. این بسته علاوه بر قرار دادن خودکار کاما، اعداد را به صورت استاندارد و مرتب چاپ می‌کند.

%%%%%%%%%%%%%%%%%%%%%%%%%%%%%%%%%%%%%%%%%%%%%%%%%%%%
% جدول‌ها
%%%%%%%%%%%%%%%%%%%%%%%%%%%%%%%%%%%%%%%%%%%%%%%%%%%%
\usepackage{booktabs}
\usepackage{array}
\usepackage{longtable}
\usepackage{multirow}
\usepackage{tabularx}